\chapter{Seizures}

\section*{Definition}
A seizure happens when there is a sudden surge of electrical activity in your brain, causing changes in your movements, behavior, sensations, or awareness.

\section*{What Happens in Your Body}
Seizures can involve convulsions, staring spells, or changes in sensation. They need medical evaluation to find the underlying cause, which may include epilepsy, head injury, infection, metabolic problems, or a brain tumor.

\section*{Causes}
\begin{itemize}
  \item Epilepsy — a condition where you have repeated, unprovoked seizures
  \item Head injury
  \item Infection of the brain (such as meningitis or encephalitis)
  \item Metabolic imbalances (such as low blood sugar or electrolyte problems)
  \item Brain tumor
\end{itemize}

\section*{When to See a Doctor}
Seek emergency care immediately if:
\begin{itemize}
  \item This is a first seizure, the seizure lasts five minutes or longer, seizures repeat without recovery, breathing is impaired, or a serious injury occurred
  \item The person is pregnant, has diabetes, or the seizure occurred in water
\end{itemize}

\section*{Related Symptoms}
\begin{itemize}
  \item Post-seizure confusion
  \item Muscle soreness
  \item Headache
  \item Amnesia of event
  \item Temporary weakness
\end{itemize}

\section*{Self-Care / Home Management}
\begin{itemize}
  \item Create a safe environment by removing sharp objects and padding hard surfaces near where seizures occur.
  \item During a seizure, do not restrain the person or put anything in their mouth. Time the seizure.
  \item After a seizure, place the person on their side to keep the airway clear.
  \item Wear a medical alert bracelet so others know about your condition.
  \item Avoid triggers: missed medication, sleep deprivation, alcohol, and flashing lights if photosensitive.
  \item Call emergency services if a seizure lasts more than 5 minutes or repeats without recovery.
\end{itemize}

\section*{How Your Doctor Will Evaluate You}
\begin{itemize}
  \item Detailed neurological history includes onset pattern, progression, triggers, associated symptoms, and functional impact.
  \item Comprehensive neurological examination assesses mental status, cranial nerves, motor/sensory function, reflexes, coordination, and gait.
  \item Neuroimaging (CT, MRI) is often indicated to evaluate for structural, vascular, or demyelinating causes.
  \item Specialized testing (EEG for seizures, nerve conduction studies for neuropathy, lumbar puncture for meningitis) is guided by clinical suspicion.
\end{itemize}

\section*{Treatment Options}
\begin{itemize}
  \item Treatment depends on the cause. A clinician may prescribe an antiseizure medicine after assessing recurrence risk, the seizure type, and the underlying condition.
  \item Prolonged or repeated seizures require emergency treatment, airway support, and medicines such as a benzodiazepine under an emergency plan.
  \item Do not stop prescribed antiseizure medicine abruptly; neurology follow-up and safety counselling are important.
\end{itemize}

\section*{References}
\begin{thebibliography}{99}
\bibitem{seizures:ilae-seizure} Fisher RS, et al. ``Operational classification of seizure types by the ILAE.'' \emph{Epilepsia}. 2024;65(3):521-535.
\bibitem{seizures:uptodate-seizure} Abou-Khalil BW, et al. Evaluation and management of first seizure in adults. In: UpToDate, Post TW (Ed), UpToDate, Waltham, MA; 2025.
\bibitem{seizures:lancet-neurol-seizure} Falco-Walter J, et al. ``The new ILAE seizure classification.'' \emph{Lancet Neurol}. 2023;22(6):520-532.
\bibitem{seizures:nejm-first-seizure} Krumholz A, et al. ``First Seizure in Adults.'' \emph{N Engl J Med}. 2023;388(15):1402-1412.
\bibitem{seizures:harrison-epilepsy} Jameson JL, et al. \emph{Harrison\'s Principles of Internal Medicine}. 21st ed. McGraw-Hill; 2022. Chapter 100: Seizures and Epilepsy.
\bibitem{seizures:nice-epilepsy} National Institute for Health and Care Excellence. ``Epilepsies: diagnosis and management.'' NICE Guideline NG217; 2023.
\end{thebibliography}
